\documentclass[10pt,a4paper]{amsart}
\usepackage{amsmath,amssymb,amsthm,graphicx}
% Retained minimal preamble for the Stacks Project chapter flat.tex (More on Flatness).
% Verbatim excerpt of the chapter preamble: Xy-pic setup (native preamble lines 14-18) and the native
% theorem declarations with the mathematical macros (native preamble lines 154-217). Copied mechanically;
% nothing added, removed or redefined. The wrapper class and the amsmath/amssymb/amsthm/graphicx base
% packages are supplied by the assembler, so the class/externaldocument part of the native preamble is omitted.
% For commutative diagrams we use Xy-pic
\usepackage[all]{xy}

% We use 2cell for 2-commutative diagrams.
\xyoption{2cell}
% Theorem environments.
%
\theoremstyle{plain}
\newtheorem{theorem}[subsection]{Theorem}
\newtheorem{proposition}[subsection]{Proposition}
\newtheorem{lemma}[subsection]{Lemma}

\theoremstyle{definition}
\newtheorem{definition}[subsection]{Definition}
\newtheorem{example}[subsection]{Example}
\newtheorem{exercise}[subsection]{Exercise}
\newtheorem{situation}[subsection]{Situation}

\theoremstyle{remark}
\newtheorem{remark}[subsection]{Remark}
\newtheorem{remarks}[subsection]{Remarks}

\numberwithin{equation}{subsection}

% Macros
%
\def\lim{\mathop{\mathrm{lim}}\nolimits}
\def\colim{\mathop{\mathrm{colim}}\nolimits}
\def\Spec{\mathop{\mathrm{Spec}}}
\def\Hom{\mathop{\mathrm{Hom}}\nolimits}
\def\Ext{\mathop{\mathrm{Ext}}\nolimits}
\def\SheafHom{\mathop{\mathcal{H}\!\mathit{om}}\nolimits}
\def\SheafExt{\mathop{\mathcal{E}\!\mathit{xt}}\nolimits}
\def\Sch{\mathit{Sch}}
\def\Mor{\mathop{\mathrm{Mor}}\nolimits}
\def\Ob{\mathop{\mathrm{Ob}}\nolimits}
\def\Sh{\mathop{\mathit{Sh}}\nolimits}
\def\NL{\mathop{N\!L}\nolimits}
\def\CH{\mathop{\mathrm{CH}}\nolimits}
\def\proetale{{pro\text{-}\acute{e}tale}}
\def\etale{{\acute{e}tale}}
\def\QCoh{\mathit{QCoh}}
\def\Ker{\mathop{\mathrm{Ker}}}
\def\Im{\mathop{\mathrm{Im}}}
\def\Coker{\mathop{\mathrm{Coker}}}
\def\Coim{\mathop{\mathrm{Coim}}}

% Boxtimes
%
\DeclareMathSymbol{\boxtimes}{\mathbin}{AMSa}{"02}

%
% Macros for moduli stacks/spaces
%
\def\QCohstack{\mathcal{QC}\!\mathit{oh}}
\def\Cohstack{\mathcal{C}\!\mathit{oh}}
\def\Spacesstack{\mathcal{S}\!\mathit{paces}}
\def\Quotfunctor{\mathrm{Quot}}
\def\Hilbfunctor{\mathrm{Hilb}}
\def\Curvesstack{\mathcal{C}\!\mathit{urves}}
\def\Polarizedstack{\mathcal{P}\!\mathit{olarized}}
\def\Complexesstack{\mathcal{C}\!\mathit{omplexes}}
% \Pic is the operator that assigns to X its picard group, usage \Pic(X)
% \Picardstack_{X/B} denotes the Picard stack of X over B
% \Picardfunctor_{X/B} denotes the Picard functor of X over B
\def\Pic{\mathop{\mathrm{Pic}}\nolimits}
\def\Picardstack{\mathcal{P}\!\mathit{ic}}
\def\Picardfunctor{\mathrm{Pic}}
\def\Deformationcategory{\mathcal{D}\!\mathit{ef}}
\usepackage{cleveref}
\crefname{originaldefinition}{original definition}{original definitions}
\Crefname{originaldefinition}{Original definition}{Original definitions}
\makeatletter
\newlabel{definition-one-step-devissage-at-x}{{05H5}{source}{}{sourcepointer-1}{}}
\newlabel{definition-one-step-devissage-at-x@cref}{{[originaldefinition][0][]05H5}{[1][0][]source}{}{}{}}
\makeatother
\crefname{originallemma}{original lemma}{original lemmas}
\Crefname{originallemma}{Original lemma}{Original lemmas}
\makeatletter
\newlabel{lemma-base-change-one-step}{{05H7}{source}{}{sourcepointer-2}{}}
\newlabel{lemma-base-change-one-step@cref}{{[originallemma][0][]05H7}{[1][0][]source}{}{}{}}
\makeatother
\title{Base change of a one step devissage at a point under a residue-field trivial pointed morphism}
\author{The Stacks Project Authors}
\date{stacks-project:flat.tex:local-snapshot-20261011; GNU FDL 1.2 or later}
\begin{document}
\maketitle
\noindent\textit{Editorial note.} All delivered mathematical text is the author's own, copied from the pinned Stacks Project chapter file flat.tex; no statement or proof step was rewritten, reordered or removed.\par\medskip
\noindent\textit{Editorial note.} The preamble is the chapter's own retained minimal preamble (sources/stacks-flat/retained-preamble.tex), a verbatim excerpt of the native Xy-pic setup and of the native theorem declarations and mathematical macros; the counted span uses no commutative diagram, but the chapter preamble is retained unedited rather than trimmed. No macro was invented or redefined.\par\medskip
\noindent\textit{Editorial note.} Both referenced results live in the same chapter, so their locators point at the \textbackslash{}label line inside the pinned chapter file, each recorded with its stable Stacks tag.\par\medskip
\noindent\textit{Allowed standard prerequisites (editorial).}\begin{itemize}
\item One step devissage of a finite type module at a point. Same chapter, Definition environment labelled definition-one-step-devissage-at-x, source line 829 (Stacks tag 05H5). Supplies the definition whose conditions (1)-(4) are checked in the second half of the counted proof; the verification itself stays inside the counted span.
\item Base change of a one step devissage over a point. Same chapter, Lemma labelled lemma-base-change-one-step, source line 879 (Stacks tag 05H7). Supplies the over-s' half of the conclusion, invoked at proof lines 942-945; the extra fibre-identification argument for the pointed version is retained in the counted span.
\end{itemize}\medskip
\ifdefined\hypertarget\hypertarget{sourcepointer-1}{}\fi
\noindent\textit{Source reference.} Original definition 05H5: sources/stacks-flat/pinned-original-source.tex, source lines 829--829. Same-chapter pointer to the definition invoked in the statement (line 927) and again at proof line 946; it names the four conditions that the counted proof transports, and the transport argument itself is retained verbatim in the delivered body.\par\medskip
\ifdefined\hypertarget\hypertarget{sourcepointer-2}{}\fi
\noindent\textit{Source reference.} Original lemma 05H7: sources/stacks-flat/pinned-original-source.tex, source lines 879--879. Same-chapter pointer used twice, in the statement at line 933 and in the proof at line 943, for the over-s' half of the conclusion; the residue-field argument that turns it into the at-a-point statement is the counted obligation and stays in the delivered body.\par\medskip
\begin{lemma}
\label{lemma-base-change-one-step-at-x}
Let $S$, $X$, $\mathcal{F}$, $x$, $s$ be as in
Definition \ref{definition-one-step-devissage-at-x}.
Let $(Z, Y, i, \pi, \mathcal{G}, z, y)$ be a one step d\'evissage
of $\mathcal{F}/X/S$ at $x$.
Let $(S', s') \to (S, s)$ be a morphism of pointed schemes
which induces an isomorphism $\kappa(s) = \kappa(s')$.
Let $(Z', Y', i', \pi', \mathcal{G}')$ be as constructed in
Lemma \ref{lemma-base-change-one-step}
and let $x' \in X'$ (resp.\ $z' \in Z'$, $y' \in Y'$) be the
unique point mapping to both $x \in X$ (resp.\ $z \in Z$, $y \in Y$)
and $s' \in S'$.
If $S'$ is affine, then $(Z', Y', i', \pi', \mathcal{G}', z', y')$
is a one step d\'evissage of $\mathcal{F}'/X'/S'$ at $x'$.
\end{lemma}

\begin{proof}
By
Lemma \ref{lemma-base-change-one-step}
$(Z', Y', i', \pi', \mathcal{G}')$ is a one step d\'evissage of
$\mathcal{F}'/X'/S'$ over $s'$. Properties (1) -- (4) of
Definition \ref{definition-one-step-devissage-at-x}
hold for $(Z', Y', i', \pi', \mathcal{G}', z', y')$
as the assumption that $\kappa(s) = \kappa(s')$ insures that the fibres
$X'_{s'}$, $Z'_{s'}$, and $Y'_{s'}$ are isomorphic to
$X_s$, $Z_s$, and $Y_s$.
\end{proof}

\end{document}
